\documentclass[times]{san}

\usepackage{graphicx}
\usepackage{float}

\addbibresource{bibliografia.bib}

\uczelnia{Społeczna Akademia Nauk}
\przedmiot{Systemy Szkieletowe}
\autor{Maura Wasielewska}
\nralbumu{121466}
\grupa{2}
\tytul{System rezerwacji sal}
\prowadzacy{mgr inż. Krystian Gumiński}
\miasto{Łódź}
\rok{2026}

\newcommand{\screen}[3]{
\begin{figure}[H]
    \centering
    \includegraphics[width=0.85\textwidth]{#1}
    \caption{#2}
    \label{#3}
\end{figure}
}



\begin{document}

\maketitlepage
\spistresci

\section{Wstęp}

W ramach projektu została opracowana aplikacja internetowa służąca do rezerwacji sal. System umożliwia użytkownikom przeglądanie dostępnych pomieszczeń, wyszukiwanie sal oraz dokonywanie rezerwacji w określonym terminie.
Aplikacja została zaprojektowana jako system składający się z części frontendowej, backendowej oraz bazy danych. W projekcie wykorzystano język Python i framework Flask, natomiast interfejs użytkownika wykonano przy użyciu HTML, CSS i JavaScript. Dane przechowywane są w bazie SQLite z wykorzystaniem SQLAlchemy.


\section{Cel projektu}

Celem projektu jest stworzenie prostego systemu umożliwiającego sprawne zarządzanie rezerwacjami sal. System pozwala użytkownikowi na rejestrację, logowanie, przeglądanie i wyszukiwanie sal, sprawdzanie ich dostępności oraz tworzenie i anulowanie rezerwacji. Dodatkowo użytkownik może korzystać z kalendarza oraz sprawdzać zajętość pomieszczeń.
Drugim istotnym celem było zastosowanie mechanizmu kontroli konfliktów terminów, który uniemożliwia zarezerwowanie tej samej sali przez dwóch użytkowników w nakładających się godzinach.System posiada również konto administratora, który może zarządzać dostępnymi salami.

\section{Przegląd istniejących rozwiązań}

Na rynku dostępnych jest wiele systemów umożliwiających rezerwację pomieszczeń. Są one wykorzystywane między innymi w przedsiębiorstwach, szkołach i uczelniach.



\subsection{Microsoft Outlook}
Microsoft Outlook umożliwia planowanie spotkań oraz rezerwowanie dostępnych zasobów, w tym sal konferencyjnych. System jest zintegrowany z kalendarzem i pozwala użytkownikom sprawdzać dostępność wybranych zasobów.
W porównaniu z projektowanym systemem Outlook oferuje znacznie szerszy zakres funkcji związanych z komunikacją i organizacją pracy.


\subsection{Microsoft Bookings}
Microsoft Bookings jest rozwiązaniem umożliwiającym zarządzanie terminami oraz rezerwacjami. System pozwala na określanie dostępnych terminów i zarządzanie rezerwacjami za pomocą interfejsu internetowego.
Projektowany system jest rozwiązaniem prostszym i skoncentrowanym wyłącznie na rezerwacji sal.


\subsection{Dedykowane systemy rezerwacji sal}
Na rynku dostępne są również dedykowane systemy przeznaczone do zarządzania przestrzenią i rezerwacją pomieszczeń. Oferują one między innymi kalendarze, zarządzanie użytkownikami oraz sprawdzanie dostępności sal.
Projektowany system realizuje podstawowe funkcje tego typu aplikacji, jednocześnie zachowując prostą strukturę i niewielką liczbę wymaganych komponentów.




\section{Opis architektury}

System został podzielony na trzy podstawowe warstwy: frontend, backend oraz bazę danych.
Frontend odpowiada za prezentację informacji oraz komunikację z użytkownikiem. Backend realizuje logikę aplikacji, obsługuje żądania HTTP i kontroluje dostęp do poszczególnych funkcji. Baza danych przechowuje informacje o użytkownikach, salach oraz rezerwacjach.


\begin{table}[h]
\centering
\caption{Główne komponenty architektury systemu}
\label{tab:lorem}
\begin{tabular}{|c|l|c|}
\hline
Lp. & Komponent          & Warstwa \\
\hline
1   & Frontend        & Prezentacja i obsługa użytkownika   \\
2   & Backend Flask      & Logika aplikacji i API \\
3   & SQLAlchemy      & Komunikacja z bazą danych     \\
3   & SQLite      & Przechowywanie danych    \\
3   & Docker     & Uruchamianie aplikacji w kontenerze    \\
\hline
\end{tabular}
\end{table}

\section{Opis implementacji}

W rozdziale przedstawiono najważniejsze elementy implementacji systemu. Opis został podzielony na frontend, backend, bazę danych oraz elementy związane z uruchamianiem, testowaniem i monitorowaniem aplikacji.

\subsection{Frontend}
Frontend aplikacji został wykonany przy wykorzystaniu HTML, CSS oraz JavaScript.
HTML odpowiada za strukturę poszczególnych stron, CSS za wygląd interfejsu, natomiast JavaScript umożliwia komunikację z backendem i dynamiczne wyświetlanie danych.
Aplikacja posiada między innymi strony logowania i rejestracji, listę sal, wyszukiwarkę, formularz rezerwacji, listę własnych rezerwacji, kalendarz, widok zajętości sal oraz panel administratora.

\screen{01-dashboard.png}
{Panel główny użytkownika systemu.}
{fig:dashboard}

\screen{02-dashboard.png}
{Panel główny użytkownika systemu.}
{fig:dashboard}

\screen{02-logowanie.png}
{Formularz logowania użytkownika.}
{fig:logowanie}

\screen{03-rejestracja.png}
{Formularz rejestracji nowego użytkownika.}
{fig:rejestracja}

\screen{04-sale.png}
{Lista dostępnych sal.}
{fig:sale}

\screen{05-wyszukiwanie.png}
{Formularz wyszukiwania i filtrowania sal.}
{fig:wyszukiwanie}


\screen{07-rezerwacja.png}
{Formularz tworzenia rezerwacji sali.}
{fig:rezerwacja}

\screen{08-konflikt.png}
{Odrzucenie rezerwacji w przypadku konfliktu terminów.}
{fig:konflikt}

\screen{09-moje-rezerwacje.png}
{Lista rezerwacji użytkownika.}
{fig:moje-rezerwacje}

\screen{10-kalendarz.png}
{Widok kalendarza rezerwacji.}
{fig:kalendarz}

\screen{10-kalendarz-konflikt.png}
{Widok kalendarza rezerwacji.}
{fig:kalendarz}



\subsection{Backend}
Backend został wykonany w języku Python z wykorzystaniem frameworka Flask.
Odpowiada on za obsługę użytkowników, sal i rezerwacji oraz komunikację z bazą danych. Backend udostępnia endpointy wykorzystujące między innymi metody HTTP GET, POST i DELETE.
Podczas tworzenia rezerwacji system sprawdza poprawność danych oraz dostępność wybranego terminu. Jeżeli istnieje konflikt z inną rezerwacją, operacja zostaje odrzucona.
Backend odpowiada również za kontrolę uprawnień. Funkcje administracyjne są dostępne wyłącznie dla użytkownika posiadającego rolę administratora.

\screen{admin.png}
{Panel administratora systemu.}
{fig:admin}

\screen{dodawanie-sali.png}
{Formularz dodawania sali przez administratora.}
{fig:dodawanie-sali}

\subsection{Baza Danych}
Do przechowywania danych wykorzystano bazę SQLite. Komunikacja pomiędzy aplikacją a bazą danych realizowana jest przy pomocy biblioteki SQLAlchemy.
W bazie przechowywane są informacje dotyczące użytkowników, sal oraz rezerwacji.
Model użytkownika zawiera między innymi identyfikator, imię, adres e-mail, hasło oraz rolę.
Model sali przechowuje nazwę sali, budynek, piętro, pojemność, wyposażenie oraz status.
Model rezerwacji zawiera informacje o użytkowniku, sali, dacie oraz godzinach rozpoczęcia i zakończenia rezerwacji.
Hasła użytkowników są przechowywane w postaci zahaszowanej.


\subsection{Docker}
Aplikacja została przygotowana do uruchamiania z wykorzystaniem technologii Docker.
W projekcie znajduje się plik Dockerfile, który określa sposób przygotowania środowiska aplikacji. Wykorzystano obraz Python 3.11 oraz plik requirements.txt zawierający wymagane biblioteki.

Obraz aplikacji został utworzony za pomocą polecenia:
docker build -t system-rezerwacji-sal .

Zastosowanie Dockera pozwala na uruchomienie aplikacji w przygotowanym i powtarzalnym środowisku.

\screen{docker.png}
{Budowanie obrazu aplikacji za pomocą technologii Docker.}
{fig:docker}

\subsection{Testy}
Do testowania aplikacji wykorzystano bibliotekę unittest.
Przygotowane testy sprawdzają między innymi dostępność strony głównej, strony logowania oraz endpointu odpowiedzialnego za listę sal.
Testy uruchamiane są za pomocą polecenia:
python -m unittest discover -s tests -v
Wszystkie trzy przygotowane testy zakończyły się poprawnie:
Ran 3 tests ... OK

\screen{testy.png}
{Wynik wykonania testów jednostkowych aplikacji.}
{fig:testy}

\subsection{Monitoring}
W celu sprawdzania stanu aplikacji przygotowano endpoint:
GET /health
Endpoint sprawdza działanie aplikacji oraz możliwość wykonania operacji na bazie danych.
W przypadku prawidłowego działania zwracany jest status ok oraz informacja, że baza danych działa poprawnie. Mechanizm umożliwia szybkie sprawdzenie podstawowego stanu systemu.

\screen{health.png}
{Wynik działania endpointu monitorującego stan aplikacji.}
{fig:health}

\section{Podsumowanie i wnioski}
W ramach projektu utworzono internetowy system rezerwacji sal umożliwiający użytkownikom wyszukiwanie pomieszczeń, sprawdzanie ich dostępności oraz tworzenie i anulowanie rezerwacji.
System posiada mechanizm zapobiegania konfliktom terminów oraz podział użytkowników na zwykłych użytkowników i administratorów. Administrator ma dodatkowo możliwość zarządzania dostępnymi salami.
Do wykonania projektu wykorzystano Python, Flask, HTML, CSS, JavaScript, SQLite, SQLAlchemy oraz Docker. Kod projektu jest zarządzany przy użyciu Git i przechowywany na platformie GitHub.
Projekt pozwolił na praktyczne wykorzystanie wiedzy związanej z tworzeniem aplikacji internetowych, projektowaniem baz danych, tworzeniem API, testowaniem oraz uruchamianiem aplikacji w środowisku kontenerowym

\spistabel
\spisrysunkow
\nocite{*}
\printbibliography

\end{document}